\section{Usage}

The recommended environment is the Conda environment \texttt{MOF\_sim}.
It should contain Python, ASE, LAMMPS, CoolProp, matplotlib, and optional OVITO support.

\subsection{Dry Run}

\begin{lstlisting}[language=bash]
conda activate MOF_sim
python benchmark.py benchmark_tao2022_mof5_co2.json --dry-run
\end{lstlisting}

This prints the resolved run plan and does not write files.

\subsection{Prepare Files}

\begin{lstlisting}[language=bash]
conda activate MOF_sim
python benchmark.py benchmark_tao2022_mof5_co2.json --prepare --new-run
\end{lstlisting}

This materializes the benchmark directory with generated LAMMPS inputs and copied source files.

\subsection{Run Smoke Test}

\begin{lstlisting}[language=bash]
conda activate MOF_sim
python benchmark.py benchmark_tao2022_mof5_co2.json --run-test --new-run
\end{lstlisting}

This runs a short GCMC test to verify that the generated LAMMPS files are readable and that insertion can occur.

\subsection{Run Full Isotherm}

\begin{lstlisting}[language=bash]
conda activate MOF_sim
python benchmark.py benchmark_tao2022_mof5_co2.json --run-isotherm --new-run --jobs 4
\end{lstlisting}

The \texttt{--jobs} option parallelizes pressure points.
It should be chosen based on available CPU resources and memory.

\subsection{Evaluate Existing Run}

\begin{lstlisting}[language=bash]
conda activate MOF_sim
python benchmark.py benchmark_tao2022_mof5_co2.json --evaluate-all --run-id <run_id>
\end{lstlisting}

This evaluates an already completed run against all active reference datasets.
